\documentclass[a4paper,11pt]{article}

% --- Sprache & Kodierung ---
\usepackage[utf8]{inputenc}
\usepackage[T1]{fontenc}
\usepackage[ngerman]{babel}
\usepackage{lmodern}
\usepackage[scaled=0.92]{helvet}          % Helvetica (praktisch identisch mit Arial)
\renewcommand{\familydefault}{\sfdefault} % serifenlose Schrift als Standard
\usepackage{microtype}

% --- Seitenlayout ---
\usepackage[a4paper,margin=2.5cm,headheight=14pt]{geometry}
\setlength{\parindent}{0pt}
\setlength{\parskip}{0.6em}

% --- Mathematik ---
\usepackage{amsmath}

% --- Tabellen ---
\usepackage{booktabs}
\usepackage{tabularx}
\usepackage{array}

% --- Farben ---
\usepackage[table]{xcolor}
\definecolor{hsblue}{HTML}{1F4E79}
\definecolor{hslight}{HTML}{E8F0F8}
\definecolor{okgreen}{HTML}{2E7D32}
\definecolor{okbg}{HTML}{E8F5E9}
\definecolor{warnred}{HTML}{C62828}
\definecolor{warnbg}{HTML}{FDECEA}
\definecolor{codebg}{HTML}{F2F2F2}

% --- Boxen, Listen, Überschriften, Kopfzeile ---
\usepackage[most]{tcolorbox}
\usepackage{enumitem}
\setlist{itemsep=0.25em, topsep=0.3em}
\usepackage{titlesec}
\titleformat{\section}{\Large\bfseries\color{hsblue}}{\thesection}{0.8em}{}[{\color{hsblue}\titlerule[0.8pt]}]
\titleformat{\subsection}{\large\bfseries\color{hsblue}}{\thesubsection}{0.8em}{}
\usepackage{fancyhdr}
\pagestyle{fancy}
\fancyhf{}
\fancyhead[L]{\small\color{gray}ADR-002}
\fancyhead[R]{\small\color{gray}Mobile UX -- Chat-App}
\fancyfoot[C]{\small\thepage}
\renewcommand{\headrulewidth}{0.4pt}

% --- Links ---
\usepackage{hyperref}
\hypersetup{
    colorlinks=true,
    linkcolor=hsblue,
    urlcolor=hsblue,
    pdftitle={ADR-002: State-Management und Datenaktualisierung in der Chat-App},
    pdfauthor={Gruppe 5}
}

% --- Eigene Befehle ---
\newcommand{\code}[1]{\colorbox{codebg}{\texttt{\small #1}}}

\newtcolorbox{entscheidung}{colback=okbg, colframe=okgreen, fonttitle=\bfseries,
    title=Entscheidung, arc=2mm, boxrule=0.8pt}
\newtcolorbox{proBox}{colback=okbg, colframe=okgreen, fonttitle=\bfseries,
    title=Positive Folgen, arc=2mm, boxrule=0.8pt}
\newtcolorbox{contraBox}{colback=warnbg, colframe=warnred, fonttitle=\bfseries,
    title=Negative Folgen \& Risiken, arc=2mm, boxrule=0.8pt}

\begin{document}

% ================= Titelblock =================
\begin{tcolorbox}[colback=hsblue, colframe=hsblue, arc=2mm, left=6mm, right=6mm, top=5mm, bottom=5mm]
    \color{white}
    {\small Architecture Decision Record}\\[0.3em]
    {\LARGE\bfseries ADR-002: State-Management und\\[0.2em] Datenaktualisierung in der Chat-App}\\[0.8em]
    {\small Gruppe 5 \quad$\cdot$\quad Mobile UX, Hochschule Esslingen \quad$\cdot$\quad \today}
\end{tcolorbox}

\vspace{0.5em}
\begin{tabularx}{\textwidth}{@{}l X@{}}
    \textbf{Status:}  & \colorbox{hsblue}{\color{white}\bfseries\,Vorgeschlagen\,} \\[0.3em]
    \textbf{Baut auf:} & ADR-001 (Wahl von React~19 mit Vite \& TypeScript) \\
\end{tabularx}

% ================= 1. Titel =================
\section{Titel und Nummer}
\textbf{ADR-002: State-Management und Datenaktualisierung in der Chat-App}

% ================= 2. Status =================
\section{Status}
\textbf{Vorgeschlagen}

% ================= 3. Kontext =================
\section{Kontext}
Mit ADR-001 wurde React~19 mit Vite und TypeScript als Frontend-Framework
festgelegt. Offen ist, \textbf{wie der Anwendungszustand verwaltet} wird und
\textbf{wie neue Daten vom Chat-Server}
(\url{https://www2.hs-esslingen.de/~chkohl/chat/}) in die Oberfläche gelangen.

\subsection*{Zu verwaltender Zustand}
\begin{itemize}
    \item \textbf{Authentifizierung:} Token nach der Registrierung/Anmeldung.
          Es muss über Seiten-Reloads erhalten bleiben und an jeden
          API-Aufruf angehängt werden.
    \item \textbf{Chat-Daten:} Nachrichten des aktiven Raums, Raumliste,
          Mitglieder, Suchergebnisse.
    \item \textbf{UI-Zustand:} Aktiver Tab (\textit{API-Test} /
          \textit{Live-Chat}), geöffnete Modals, Formulareingaben, Lade- und
          Fehlerzustände.
    \item \textbf{API-Testbench:} Fortschritt des 4-stufigen Ablaufs,
          Zeitstempel für die 1s-Delay-Validierung, empfangene
          HTTP-Fehlercodes (451--469).
\end{itemize}

\subsection*{Randbedingungen}
\begin{itemize}
    \item Der Chat-Server bietet eine reine HTTP-Schnittstelle (GET mit
          Query-Parametern / POST mit JSON-Body). Neue Nachrichten müssen
          daher \textbf{per Polling} abgefragt werden.
    \item Der Zustand wird nur von wenigen Komponenten gemeinsam genutzt
          (Header, Chat-Seite, Modals).
    \item Der Code soll übersichtlich bleiben und ohne lange Einarbeitung
          nachvollziehbar sein (vgl.\ ADR-001).
\end{itemize}

% ================= 4. Optionen =================
\section{Optionen und Bewertung}
Es wurden drei Optionen für das State-Management betrachtet:
\begin{enumerate}
    \item \textbf{Option A: React-Bordmittel} (\code{useState}, \code{useReducer},
          \code{useRef}, Context API), gekapselt in Custom Hooks
          (\code{useAuth}, \code{useChat})
    \item \textbf{Option B: Zustand} (schlanke externe State-Bibliothek)
    \item \textbf{Option C: Redux Toolkit mit RTK Query}
\end{enumerate}

\subsection*{Gewichtete Entscheidungsmatrix}
Bewertungsskala: 1 (schlecht) bis 5 (sehr gut). Der Gesamtwert ergibt sich
als Summe aus \textit{Gewichtung $\times$ Punktzahl}. Die gewählte Option ist
grün hinterlegt.

\begin{table}[h]
\centering
\small
\renewcommand{\arraystretch}{1.35}
\begin{tabularx}{\textwidth}{>{\raggedright\arraybackslash}X c >{\columncolor{okbg}}c c c}
\toprule
\rowcolor{hslight}
\textbf{Kriterium} & \textbf{Gewichtung} & \textbf{A: React} & \textbf{B: Zustand} & \textbf{C: Redux} \\
\midrule
Geringer Boilerplate \& Übersichtlichkeit & 25\,\% & 5 & 4 & 2 \\
Eignung für Polling \& Token-Handling     & 25\,\% & 4 & 4 & 5 \\
TypeScript-Integration                    & 20\,\% & 5 & 5 & 4 \\
Keine zusätzlichen Abhängigkeiten         & 15\,\% & 5 & 3 & 2 \\
Skalierbarkeit \& Debugging-Tools         & 15\,\% & 3 & 4 & 5 \\
\midrule
\textbf{Gesamtbewertung} & \textbf{100\,\%} & \textbf{4{,}45} & \textbf{4{,}05} & \textbf{3{,}60} \\
\bottomrule
\end{tabularx}
\caption{Gewichtete Entscheidungsmatrix (Punktzahlen 1--5)}
\end{table}

\textit{Rechenweg:}
\begin{align*}
\text{A:}\quad & 0{,}25 \cdot 5 + 0{,}25 \cdot 4 + 0{,}20 \cdot 5 + 0{,}15 \cdot 5 + 0{,}15 \cdot 3 = \mathbf{4{,}45}\\
\text{B:}\quad & 0{,}25 \cdot 4 + 0{,}25 \cdot 4 + 0{,}20 \cdot 5 + 0{,}15 \cdot 3 + 0{,}15 \cdot 4 = \mathbf{4{,}05}\\
\text{C:}\quad & 0{,}25 \cdot 2 + 0{,}25 \cdot 5 + 0{,}20 \cdot 4 + 0{,}15 \cdot 2 + 0{,}15 \cdot 5 = \mathbf{3{,}60}
\end{align*}

\subsection*{Detaillierte Bewertung der Alternativen}
\begin{itemize}
    \item \textbf{Redux Toolkit mit RTK Query (Option C):}
    \begin{itemize}
        \item \textit{Vorteile:} Eingebautes Polling
              (\code{pollingInterval}), Caching, sehr gute DevTools,
              bewährt in großen Projekten.
        \item \textit{Nachteile:} Hoher Einarbeitungs- und
              Konfigurationsaufwand (Store, Slices, API-Definitionen) -- für
              den Umfang der Übung überdimensioniert.
    \end{itemize}
    \item \textbf{Zustand (Option B):}
    \begin{itemize}
        \item \textit{Vorteile:} Sehr wenig Code, globaler Store ohne
              Provider, gute TypeScript-Unterstützung,
              \code{persist}-Middleware für das Token.
        \item \textit{Nachteile:} Zusätzliche Abhängigkeit. Bei wenigen
              gemeinsam genutzten Zuständen ist der Vorteil gegenüber
              React-Bordmitteln gering.
    \end{itemize}
    \item \textbf{React-Bordmittel in Custom Hooks (Option A):}
    \begin{itemize}
        \item \textit{Vorteile:} Keine zusätzliche Bibliothek. Die Logik ist
              sauber in \code{useAuth} und \code{useChat} gekapselt. Polling
              über \code{useEffect} + \code{setInterval}, mit \code{useRef}
              für den letzten Nachrichtenstand, ist gut kontrollierbar.
        \item \textit{Nachteile:} Kein eingebautes Caching. Bei stark
              wachsendem Zustand kann Context unnötige Re-Renders verursachen.
    \end{itemize}
\end{itemize}

% ================= 5. Entscheidung =================
\section{Entscheidung}
\begin{entscheidung}
Wir entscheiden uns für \textbf{Option A: React-Bordmittel, gekapselt in
Custom Hooks}.
\end{entscheidung}

Umsetzung im Detail:
\begin{itemize}
    \item \textbf{\code{useAuth} (via \code{AuthContext}):} Hält das
          Token, speichert es in \code{localStorage} und stellt
          \code{login}, \code{register} und \code{logout} bereit. Der
          Context wird in der App-Wurzel bereitgestellt, da Header, Chat und
          API-Test das Token benötigen.
    \item \textbf{\code{useChat}:} Verwaltet Nachrichten, Räume und
          Mitglieder des aktiven Raums. Neue Nachrichten werden per
          \textbf{Polling} abgefragt (\code{setInterval} in
          \code{useEffect}, Aufräumen im Cleanup). Über \code{useRef}
          wird die ID bzw.\ der Zeitstempel der letzten Nachricht gemerkt,
          sodass nur neue Nachrichten angehängt werden.
    \item \textbf{Lokaler UI-Zustand} (Modals, Formulareingaben, aktiver Tab)
          bleibt per \code{useState} in der jeweiligen Komponente.
    \item \textbf{\code{src/api/}} bleibt zustandslos. Die Hooks rufen die
          API-Funktionen auf und übergeben das Token explizit.
\end{itemize}

% ================= 6. Konsequenzen =================
\section{Konsequenzen}
\begin{proBox}
\begin{itemize}[leftmargin=*]
    \item \textbf{Keine zusätzlichen Abhängigkeiten:} Das Projekt bleibt
          schlank, der Build schnell.
    \item \textbf{Klare Verantwortlichkeiten:} Komponenten zeigen nur an,
          Logik und Zustand liegen in den Hooks -- passend zur
          Schichtenarchitektur aus ADR-001.
    \item \textbf{Nachvollziehbares Polling:} Start, Stopp und Intervall sind
          explizit im Code sichtbar. Beim Raumwechsel oder Logout wird das
          Intervall sauber beendet.
    \item \textbf{Persistente Anmeldung:} Durch \code{localStorage} bleibt
          das Token über Seiten-Reloads erhalten.
\end{itemize}
\end{proBox}

\begin{contraBox}
\begin{itemize}[leftmargin=*]
    \item \textbf{Polling-Last:} Regelmäßige Anfragen erzeugen
          Netzwerkverkehr, auch wenn keine neuen Nachrichten vorliegen. Das
          Intervall muss passend gewählt werden (z.\,B. 2--3\,s).
    \item \textbf{Fehleranfälligkeit bei \code{useEffect}:} Fehlendes
          Cleanup führt zu doppelten Intervallen (insbesondere im
          React-StrictMode) oder zu veralteten Werten in Closures.
    \item \textbf{Token in \code{localStorage}:} Anfällig bei XSS. Für die
          Übung akzeptabel, im Produktivbetrieb wären HttpOnly-Cookies
          vorzuziehen.
    \item \textbf{Migrationspfad:} Wächst der Zustand deutlich, ist ein
          Wechsel zu Zustand (Option B) mit geringem Aufwand möglich, da die
          Logik bereits in Hooks gekapselt ist.
\end{itemize}
\end{contraBox}

\end{document}
